\subsection{\texorpdfstring{同态 \(\pi_{1}(f,a)\) 的性质}{同态 pi\_1(f,a) 的性质}}

设 A 与 B 是拓扑空间，\((\mathrm{E},p)\) 是 A 的一个覆叠，\((\mathrm{E}^{\prime},p^{\prime})\) 是 B 的一个覆叠，\(f\colon A\to B\) 与 \(g\colon E\to E^{\prime}\) 是满足 \(p^{\prime}\circ g=f\circ p\) 的连续映射。设 a 是 A 的一点，并设 \(b=f(a)\)。对所有 \(\gamma\in\pi_{1}(\mathrm{A},a)\) 及每个 \(x\in E_{a}\)，依 III, p. 304 的关系式 (3)，我们有 \(g(x\cdot\gamma)=g(x)\cdot f_{*}(\gamma)\)。若图表

\[
\begin{array}{c} \mathrm{E} \xrightarrow {g} \mathrm{E} ^ {\prime} \\ \Big \downarrow p \\ \mathrm{A} \xrightarrow {f} \mathrm{B} \end{array}
\]

是笛卡尔方块，则映射 g 诱导一个从 \(E_{a}\) 到 \(E'_{b}\) 上的双射。于是 \(\pi_{1}(A,a)\) 在纤维 \(E_{a}\) 上的作用，是群 \(\pi_{1}(B,b)\) 在 \(E'_{b}\) 上的作用与同态 \(\pi_{1}(f,a)\)（从 \(\pi_{1}(A,a)\) 到 \(\pi_{1}(B,b)\) 中）的复合。

命题 1. --- 设 A 是可解圈拓扑空间，B 是局部道路连通拓扑空间，并设 \(f\colon\mathrm{A}\to\mathrm{B}\) 是连续映射。设 a 是 A 的一点，\(b=f(a)\)。假设 A 的每个覆叠都同构于 B 的某个覆叠在 \(f\) 下的逆像。那么同态 \(\pi_{1}(f,a)\colon\pi_{1}(\mathrm{A},a)\to\pi_{1}(\mathrm{B},b)\) 是单射的。

由于空间 A 是可解圈的，存在 A 的一个覆叠 E，其纤维 \(E_{a}\) 是主齐性 \(\pi_{1}(A,a)\) -集合（IV, p. 342, 定理 1）。依假设，这个作用经由同态 \(\pi_{1}(f,a)\) 由 \(\pi_{1}(B,b)\) 在 \(E_{a}\) 上的一个作用诱导。这蕴含同态 \(\pi_{1}(f,a)\) 的单射性。

命题 2. --- 设 A 是连通且局部道路连通的拓扑空间，B 是可解圈拓扑空间，并设 \(f\colon\mathrm{A}\to\mathrm{B}\) 是连续映射。设 a 是 A 的一点，\(b=f(a)\)。下列断言等价：

\begin{enumerate}
\def\labelenumi{(\roman{enumi})}
\item
  同态 \(\pi_1(f,a)\colon \pi_1(A,a)\to \pi_1(B,b)\) 是满射的。
\item
  对 B 的覆叠的每个偶 (E, E')，映射
\end{enumerate}

\[
f ^ {*} \colon \mathscr {C} _ {\mathrm{B}} (\mathrm{E}; \mathrm{E} ^ {\prime}) \rightarrow \mathscr {C} _ {\mathrm{A}} (\mathrm{A} \times_ {\mathrm{B}} \mathrm{E}; \mathrm{A} \times_ {\mathrm{B}} \mathrm{E} ^ {\prime})
\]

它把 B-态射 \(g\colon \mathrm{E}\to \mathrm{E}^{\prime}\) 对应到 A-态射

\[
f ^ {*} (g) \colon \mathrm{A} \times_ {\mathrm{B}} \mathrm{E} \rightarrow \mathrm{A} \times_ {\mathrm{B}} \mathrm{E} ^ {\prime}, \quad (x, y) \mapsto (x, g (y))
\]

是双射的。

借助同态 \(\pi_1(f, a)\)，每个 \(\pi_1(B, b)\) -集合都赋有一个 \(\pi_1(A, a)\) -集合的结构。于是条件 (i) 等价于下列条件 (i')：

(i') 对每个偶 \((\mathrm{F},\mathrm{F}^{\prime})\)（\(\pi_{1}(\mathrm{B},b)\) -集合的偶），每个 \(\pi_{1}(\mathrm{A},a)\) -态射（从 F 到 \(F^{\prime}\)）都是 \(\pi_{1}(\mathrm{B},b)\) -态射。事实上，若同态 \(\pi_{1}(f,a)\) 是满射的，则每个 \(\pi_{1}(\mathrm{A},a)\) -态射（\(\pi_{1}(\mathrm{B},b)\) -集合的态射）都是 \(\pi_{1}(\mathrm{B},b)\) -态射。反之，取一个退缩为一点的 \(\pi_{1}(\mathrm{B},b)\) -集合 F，并置 \(\mathrm{F}^{\prime}=\pi_{1}(\mathrm{B},b)/f_{*}(\pi_{1}(\mathrm{A},a))\)。从 F 到 \(F^{\prime}\) 的、以 \(f_{*}(\pi_{1}(\mathrm{A},a))\) 为像的映射是一个 \(\pi_{1}(\mathrm{A},a)\) -态射，但不是 \(\pi_{1}(\mathrm{B},b)\) -态射，如果 \(F^{\prime}\) 不退缩为一点，也就是说，如果 \(\pi_{1}(f,a)\) 不是满射的。

由于空间 B 是可解圈的，每个 \(\pi_{1}(B,b)\) -集合都同构于 \(\pi_{1}(B,b)\) -集合 \(E_{b}\)，其中 E 是 B 的覆叠（IV, p. 344, 注 1）。于是 (i') 与 (ii) 的等价性由 III, p. 310 的命题 2 得出。

命题 3. --- 设 B 是可解圈拓扑空间，A 是 B 的连通且局部道路连通子空间（例如一个连通开子集）。设 a 是 A 的一点。下列性质等价：

\begin{enumerate}
\def\labelenumi{(\roman{enumi})}
\item
  B 的每个覆叠都在 A 上可平凡化。
\item
  由典范单射导出的、从 \(\pi_{1}(\mathrm{A},a)\) 到 \(\pi_{1}(\mathrm{B},a)\) 中的同态的像退缩为单位元。
\end{enumerate}

蕴涵关系 (ii)⇒(i) 由推论 3（III, p. 309）得出。

反之，假设 B 的每个覆叠都在 A 上可平凡化。对道路单连通覆叠 \((\lambda_{a}(\mathrm{B}),\varepsilon_{\mathrm{B}})\)（IV, p. 342, 定理 1）而言尤其如此，从而同态 \(\pi_{1}(\mathrm{A},a)\to\pi_{1}(\mathrm{B},a)\) 平凡（III, p. 309, 命题 1 的推论 3）。

